% Options for packages loaded elsewhere
\PassOptionsToPackage{unicode}{hyperref}
\PassOptionsToPackage{hyphens}{url}
\documentclass[
  11pt,
]{article}
\usepackage{xcolor}
\usepackage[margin=1in]{geometry}
\usepackage{amsmath,amssymb}
\setcounter{secnumdepth}{5}
\usepackage{iftex}
\ifPDFTeX
  \usepackage[T1]{fontenc}
  \usepackage[utf8]{inputenc}
  \usepackage{textcomp} % provide euro and other symbols
\else % if luatex or xetex
  \usepackage{unicode-math} % this also loads fontspec
  \defaultfontfeatures{Scale=MatchLowercase}
  \defaultfontfeatures[\rmfamily]{Ligatures=TeX,Scale=1}
\fi
\usepackage{lmodern}
\ifPDFTeX\else
  % xetex/luatex font selection
\fi
% Use upquote if available, for straight quotes in verbatim environments
\IfFileExists{upquote.sty}{\usepackage{upquote}}{}
\IfFileExists{microtype.sty}{% use microtype if available
  \usepackage[]{microtype}
  \UseMicrotypeSet[protrusion]{basicmath} % disable protrusion for tt fonts
}{}
\makeatletter
\@ifundefined{KOMAClassName}{% if non-KOMA class
  \IfFileExists{parskip.sty}{%
    \usepackage{parskip}
  }{% else
    \setlength{\parindent}{0pt}
    \setlength{\parskip}{6pt plus 2pt minus 1pt}}
}{% if KOMA class
  \KOMAoptions{parskip=half}}
\makeatother
\setlength{\emergencystretch}{3em} % prevent overfull lines
\providecommand{\tightlist}{%
  \setlength{\itemsep}{0pt}\setlength{\parskip}{0pt}}
\usepackage{needspace}
\usepackage{bookmark}
\IfFileExists{xurl.sty}{\usepackage{xurl}}{} % add URL line breaks if available
\urlstyle{same}
\hypersetup{
  pdftitle={Selective Tests of Penalized-Fit Shapley Contributions after Cross-Validation},
  hidelinks,
  pdfcreator={LaTeX via pandoc}}

\title{Selective Tests of Penalized-Fit Shapley Contributions after
Cross-Validation}
\author{}
\date{}

\begin{document}
\maketitle

\setcounter{section}{1}

\section{Problem Setup}\label{problem-setup}

\subsection{Statistical model and fixed auxiliary
information}\label{statistical-model-and-fixed-auxiliary-information}

Consider a regression dataset \((X,\boldsymbol y)\) with \(n\)
observations and \(m\) candidate features, where
\(X=[\boldsymbol x_1,\ldots,\boldsymbol x_m]\in\mathbb R^{n\times m}\)
is fixed, \(X_{ij}\) is its \((i,j)\) entry, and
\(\boldsymbol y\in\mathbb R^n\) is a realization of

\[
\boldsymbol Y=\boldsymbol\mu(X)+\boldsymbol\varepsilon,\qquad
\boldsymbol\varepsilon\sim \mathcal N(\boldsymbol0,\sigma^2 I_n),\qquad \sigma>0\text{ known}.
\tag{1}
\]

Here \(I_n\) is the \(n\times n\) identity matrix and \(\boldsymbol0\)
denotes a zero vector of the required dimension. Write
\(\boldsymbol\mu=\boldsymbol\mu(X)\) and denote the observed response by
\(\boldsymbol y_{\rm obs}\); \(\boldsymbol y\) denotes a generic
response vector. We impose no parametric representation of
\(\boldsymbol\mu\): the regression basis below defines the contribution
being tested, rather than a requirement that the true mean belong to
that basis. All probability statements condition on \(X\). They also
condition on auxiliary random choices, denoted collectively by
\(\omega\), that determine basis centers and cross-validation folds.
These choices are independent of the response noise given \(X\). The
basis fraction, candidate penalty grid, and number \(k\) of selected
features are specified before inspecting \(\boldsymbol Y\).

\subsection{Coalition models and response-dependent
tuning}\label{coalition-models-and-response-dependent-tuning}

Let \([m]=\{1,\ldots,m\}\), and fix \(k\in\{1,\ldots,m\}\). For a
prespecified fraction \(\rho\in(0,1]\), sample
\(m_{\mathrm{RBF}}=\max\{1,\lceil\rho n\rceil\}\) distinct row indices
\(i_1,\ldots,i_{m_{\mathrm{RBF}}}\) using \(\omega\), sharing these
indices across features. For feature \(j\in[m]\), define centers
\(c_{jq}=X_{i_qj}\) and, for \(q=1,\ldots,m_{\mathrm{RBF}}\),
\(i=1,\ldots,n\), and \(t\in\mathbb R\), set

\[
g_{jq}(t)=\exp\!\left[-\frac{(t-c_{jq})^2}{2s_j^2}\right],\qquad
(G_j)_{iq}=g_{jq}(X_{ij})-\frac1n\sum_{a=1}^n g_{jq}(X_{aj}).
\tag{2}
\]

The width \(s_j\) is the median of \(|X_{aj}-X_{bj}|\) over
\(1\le a<b\le n\). If this is zero, use the median strictly positive
distance; if there are no positive distances, use \(s_j=1\). Let
\(\mathbf1_n\) denote the vector of \(n\) ones. Thus
\(G_j\in\mathbb R^{n\times m_{\mathrm{RBF}}}\) is determined by \(X\)
and \(\omega\), and \(\mathbf1_n^\top G_j=\boldsymbol0^\top\). For
\(S\subseteq[m]\), let \(G_S\in\mathbb R^{n\times m_{\mathrm{RBF}}|S|}\)
concatenate the blocks indexed by \(S\), with \(G_\varnothing\) having
zero columns. These are additive univariate Gaussian basis models; the
construction contains no multivariate interaction basis.

\Needspace{7\baselineskip}

For a fixed \(\lambda>0\), define the minimized penalized training
objective

\[
L_\lambda(S;\boldsymbol y)=\min_{\beta_0,\boldsymbol\beta_{-0}}\left\{
\|\boldsymbol y-\beta_0\mathbf1_n-G_S\boldsymbol\beta_{-0}\|_2^2+n\lambda\|\boldsymbol\beta_{-0}\|_2^2\right\}.
\tag{3}
\]

Here \(\beta_0\in\mathbb R\) is the unpenalized intercept,
\(\boldsymbol\beta_{-0}\in\mathbb R^{m_{\mathrm{RBF}}|S|}\) contains the
basis coefficients. With \(P_0=\mathbf1_n\mathbf1_n^\top/n\),
\(C_0=I_n-P_0\), and \(K_S=G_SG_S^\top\), set

\[
H_S^\lambda=K_S(K_S+n\lambda I_n)^{-1},\qquad H_\varnothing^\lambda=0.
\]

For \(j\in[m]\setminus S\), centering gives

\[
\begin{aligned}
L_\lambda(S;\boldsymbol y)&=\boldsymbol y^\top(C_0-H_S^\lambda)\boldsymbol y,\\
\Delta_j^\lambda(S;\boldsymbol y)&=L_\lambda(S;\boldsymbol y)-L_\lambda(S\cup\{j\};\boldsymbol y)\\
&=\boldsymbol y^\top(H_{S\cup\{j\}}^\lambda-H_S^\lambda)\boldsymbol y.
\end{aligned}
\tag{4}
\]

The increment matrix is positive semidefinite. Indeed, adding a block
increases \(K_S\) in the Loewner order and
\(H_S^\lambda=I_n-n\lambda(K_S+n\lambda I_n)^{-1}\). Consequently
\(\Delta_j^\lambda(S;\boldsymbol y)\ge0\). Equation (4) concerns the
\emph{minimized penalized objective}, including the ridge penalty.
Replacing it with the residual sum of squares evaluated at a ridge fit
would define a different score.

We choose \(\lambda\) by five-fold cross-validation on the full-feature
model over the ordered grid

\[
\Lambda=(\lambda_1,\ldots,\lambda_9)=(10^{-6},10^{-5},\ldots,10^2).
\]

Within each fold, centers, widths, and centering means are constructed
from the training rows alone, and the penalty in (3) uses that fold's
training sample size. With \(X\) and \(\omega\) fixed, each validation
prediction is linear in \(\boldsymbol y\). Therefore there are fixed
symmetric positive semidefinite matrices \(C_\ell\) such that the pooled
validation mean squared error is

\[
\operatorname{CV}_\ell(\boldsymbol y)=\boldsymbol y^\top C_\ell \boldsymbol y,\qquad
\widehat\ell(\boldsymbol y)=\underset{\ell\in\{1,\ldots,9\}}{\operatorname{argmin}}\;\boldsymbol y^\top C_\ell \boldsymbol y,
\qquad \widehat\lambda(\boldsymbol y)=\lambda_{\widehat\ell(\boldsymbol y)}.
\tag{5}
\]

Ties are resolved toward the smaller penalty. The full-data basis is
then used with the same selected penalty for every coalition. There is
no coalition-specific tuning.

\subsection{Shapley ranking and the selected
target}\label{shapley-ranking-and-the-selected-target}

For a nonempty feature universe \(U\subseteq[m]\) and \(j\in U\), define

\[
w_{j,U}(S)=\frac{|S|!(|U|-|S|-1)!}{|U|!}
=\frac{1}{|U|\binom{|U|-1}{|S|}},
\qquad S\subseteq U\setminus\{j\},
\]

\[
A_{j,U}^\lambda
=\sum_{S\subseteq U\setminus\{j\}}w_{j,U}(S)
\bigl(H_{S\cup\{j\}}^\lambda-H_S^\lambda\bigr),
\qquad
T_{j,U}^\lambda(\boldsymbol y)=\boldsymbol y^\top A_{j,U}^\lambda \boldsymbol y.
\tag{6}
\]

\Needspace{7\baselineskip}

Thus \(T_{j,U}^\lambda\) averages the improvement from adding feature
\(j\) over all insertion orders within \(U\). The matrices are positive
semidefinite, and the allocation satisfies

\[
\sum_{j\in U}T_{j,U}^\lambda(\boldsymbol y)
=L_\lambda(\varnothing;\boldsymbol y)-L_\lambda(U;\boldsymbol y).
\tag{7}
\]

The coalition sums in (6) are evaluated exhaustively. Rank the features
by decreasing
\(T_{j,[m]}^{\widehat\lambda(\boldsymbol y)}(\boldsymbol y)\), breaking
ties by increasing feature index, and denote the unordered top-\(k\) set
by \(M(\boldsymbol y)\). Independently generate
\(V\sim\operatorname{Unif}[0,1)\) and set

\[
J(\boldsymbol y,V)=M(\boldsymbol y)_{(1+\lfloor kV\rfloor)},
\tag{8}
\]

where the elements of \(M(\boldsymbol y)\) are listed in increasing
index order. Conditional on \(M(\boldsymbol y)\), this chooses one
target uniformly. The realized value \(v\) is held fixed during
inference.

Let \(\ell_0=\widehat\ell(\boldsymbol y_{\rm obs})\),
\(\lambda_0=\lambda_{\ell_0}\), \(M_0=M(\boldsymbol y_{\rm obs})\), and
\(j=J(\boldsymbol y_{\rm obs},v)\). We consider two distinct targets for
this same selected feature:

\[
\begin{aligned}
U&=M_0 &&\text{(Selected-Features Shapley contribution)},\\
U&=[m] &&\text{(All-Features Shapley contribution)}.
\end{aligned}
\tag{9}
\]

Both use the ranking based on the full feature universe. The
Selected-Features target does not rerank features using a
restricted-universe score. The two procedures share the CV winner,
selected set, and target feature.

\subsection{Hypotheses and their
interpretation}\label{hypotheses-and-their-interpretation}

For the realized \(U,j,\ell_0\), abbreviate \(A=A_{j,U}^{\lambda_0}\)
and define the mean-response contribution

\[
T_{j,U}^{\lambda_0}(\boldsymbol\mu)=\boldsymbol\mu^\top A\boldsymbol\mu.
\]

The ideal zero-contribution hypotheses are

\[
H_{0,j,U}:T_{j,U}^{\lambda_0}(\boldsymbol\mu)=0
\quad\text{versus}\quad
H_{1,j,U}:T_{j,U}^{\lambda_0}(\boldsymbol\mu)>0.
\tag{10}
\]

These parameters refer to applying the coalition objective to the
noise-free mean vector. They are not the expectation of the observed
attribution: for fixed \(A\),
\(\mathbb E_{\boldsymbol\mu}[\boldsymbol Y^\top A\boldsymbol Y]=\boldsymbol\mu^\top A\boldsymbol\mu+\sigma^2\operatorname{tr}(A)\).

Let \(W_j=\operatorname{range}(A)\), \(P_{W_j}\) be its orthogonal
projector, and \(d_j=\dim W_j\); the dependence of \(W_j\) and \(d_j\)
on \(U\) and \(\lambda_0\) is suppressed. Since \(A\succeq0\),

\[
\boldsymbol\mu^\top A\boldsymbol\mu=0\quad\Longleftrightarrow\quad P_{W_j}\boldsymbol\mu=\boldsymbol0.
\tag{11}
\]

Moreover, positivity of every Shapley weight implies that the null in
(10) holds exactly when \(\Delta_j^{\lambda_0}(S;\boldsymbol\mu)=0\) for
every \(S\subseteq U\setminus\{j\}\). Hence the test concerns the
absence of contribution across \emph{all} coalitions in the chosen
universe. It does not isolate a coefficient in the full additive model
or assert conditional independence or causality.

The feature label need not imply a small test space. Writing
\(\mathcal V_U=\operatorname{col}(G_U)\), one has
\(W_j\subseteq\mathcal V_U\). If \(W_j=\mathcal V_U\), (10) becomes the
absence of mean signal in that entire centered basis space. Different
features can have different attribution scores but the same
zero-contribution null.

\textbf{Numerical specification of the test space.} For numerical
computation, we first compress \([G_1,\ldots,G_m]\) by retaining
singular values exceeding \(\tau\max(1,s_{\max})\), with
\(\tau=10^{-10}\), where \(s_{\max}\) is the largest singular value.
Coalition matrices are constructed in this retained space. From the
resulting target matrix, we retain eigenvectors whose eigenvalues exceed
\(\tau\max(1,a_{\max})\). Here \(a_{\max}\) is the largest eigenvalue of
the compressed target matrix. If \(Q\) contains the retained orthonormal
left singular vectors of the full basis and \(B\) contains these
orthonormal target eigenvectors in \(Q\) coordinates, the test uses

\[
W_{j,\tau}=\operatorname{col}(QB),\qquad
H_{0,j,U}^{(\tau)}:P_{W_{j,\tau}}\boldsymbol\mu=\boldsymbol0.
\tag{12}
\]

All subsequent formulas apply with this numerically specified
\(W_{j,\tau}\) as well as with ideal \(W_j\). Equivalence to the ideal
quadratic null requires agreement of the retained and ideal ranges; a
discarded small positive eigenvalue is not an exact mathematical zero.
The tolerance is therefore part of the target specification.

\subsection{Conditional inferential
procedure}\label{conditional-inferential-procedure}

For either specified test space, let \(P\) denote its orthogonal
projector (\(P_{W_j}\) or \(P_{W_{j,\tau}}\)) and
\(d_j=\operatorname{rank}(P)\) its rank. When \(d_j\ge1\), define

\[
R_j(\boldsymbol Y)=\frac{\|P\boldsymbol Y\|_2}{\sigma},\qquad
\mathcal Q_{\boldsymbol Y}=(I_n-P)\boldsymbol Y,\qquad
\boldsymbol D_j(\boldsymbol Y)=\frac{P\boldsymbol Y}{\|P\boldsymbol Y\|_2}.
\tag{13}
\]

For fixed \(P\), under \(P\boldsymbol\mu=\boldsymbol0\), the radius has
distribution \(\chi_{d_j}\) and is independent of the direction
\(\boldsymbol D_j(\boldsymbol Y)\) and orthogonal component
\(\mathcal Q_{\boldsymbol Y}\). Conditional on the realized CV winner,
selected set, and \(v\), the projector is fixed. For
\(\boldsymbol z_j=\mathcal Q_{\boldsymbol y_{\rm obs}}\) and
\(\boldsymbol u_j=\boldsymbol D_j(\boldsymbol y_{\rm obs})\), define the
response ray
\(\boldsymbol y_j(r)=\boldsymbol z_j+\sigma r\boldsymbol u_j\),
\(r\ge0\), and the event-preserving radial region

\[
\mathcal Z=
\{r\ge0:\widehat\ell(\boldsymbol y_j(r))=\ell_0,\ M(\boldsymbol y_j(r))=M_0\}.
\tag{14}
\]

The event fixes the unordered selected set, not its internal ranking.
Because \(v\) is fixed, it also preserves the target identity. For each
\(a\in M_0\) and \(b\in[m]\setminus M_0\), membership requires

\[
\boldsymbol y_j(r)^\top\bigl(A_{a,[m]}^{\lambda_0}-A_{b,[m]}^{\lambda_0}\bigr)\boldsymbol y_j(r)\ge0,
\tag{15}
\]

and, for every \(\ell\in\{1,\ldots,9\}\setminus\{\ell_0\}\), it requires

\[
\boldsymbol y_j(r)^\top(C_\ell-C_{\ell_0})\boldsymbol y_j(r)\ge0.
\tag{16}
\]

\Needspace{7\baselineskip}

Equality is admitted only when the relevant index tie rule favors the
observed choice. For either comparison, let
\(F\in\mathbb R^{n\times n}\) denote the corresponding symmetric matrix
difference in (15) or (16). The comparison is quadratic in \(r\):

\[
\sigma^2(\boldsymbol u_j^\top F\boldsymbol u_j)r^2+2\sigma(\boldsymbol u_j^\top F\boldsymbol z_j)r+\boldsymbol z_j^\top F\boldsymbol z_j\ge0.
\tag{17}
\]

Intersecting their solution sets gives a finite union of intervals, up
to endpoints of zero continuous probability. Both identically tied
comparisons and their tie rules must be retained. Importantly, the CV
matrices in (16) act on the original response space; CV need not depend
only on its projection into the full-data basis space.

Write \(f_{\chi_{d_j}}\) for the density of the chi distribution with
\(d_j\) degrees of freedom and \(\mathbf1\{\cdot\}\) for the indicator
function. Conditioning on \(\ell_0,M_0,\boldsymbol z_j,\boldsymbol u_j\)
below means conditioning on \(\widehat\ell(\boldsymbol Y)=\ell_0\),
\(M(\boldsymbol Y)=M_0\),
\(\mathcal Q_{\boldsymbol Y}=\boldsymbol z_j\), and
\(\boldsymbol D_j(\boldsymbol Y)=\boldsymbol u_j\). For positive
selection probability, the conditional density of \(R_j(\boldsymbol Y)\)
under the applicable null \(H_0\) is

\[
f(r\mid X,\omega,v,\ell_0,M_0,\boldsymbol z_j,\boldsymbol u_j;H_0)
=\frac{f_{\chi_{d_j}}(r)\mathbf1\{r\in\mathcal Z\}}
{\int_{\mathcal Z} f_{\chi_{d_j}}(t)\,dt}.
\tag{18}
\]

With \(r_{\rm obs}=R_j(\boldsymbol y_{\rm obs})\), compute the selective
upper-tail probability

\[
p_{\mathrm{selective},j,U}
=\frac{\int_{\mathcal Z\cap[r_{\rm obs},\infty)} f_{\chi_{d_j}}(t)\,dt}
{\int_{\mathcal Z} f_{\chi_{d_j}}(t)\,dt}.
\tag{19}
\]

We evaluate interval masses using chi distribution functions and
logarithmic arithmetic. For a zero-rank target, set
\(p_{\mathrm{selective},j,U}=1\). For positive-rank targets, assume
\(r_{\rm obs}>0\) and positive finite selection mass; these conditions
hold almost surely under the nondegenerate conditional model.

In exact arithmetic, (18)-(19) yield a uniform conditional \(p\)-value
under the corresponding fixed-label null, for nondegenerate conditioning
values. In particular, for every \(\boldsymbol\mu\) satisfying that null
and every \(\alpha\in(0,1)\),

\[
\begin{aligned}
\mathbb P_{\boldsymbol\mu}\!\bigl(&p_{\mathrm{selective},j,U}\le\alpha\mid
X,\omega,v,\widehat\ell(\boldsymbol Y)=\ell_0,\ M(\boldsymbol Y)=M_0,\\
&\mathcal Q_{\boldsymbol Y}=\boldsymbol z_j,\ \boldsymbol D_j(\boldsymbol Y)=\boldsymbol u_j\bigr)\le\alpha.
\end{aligned}
\tag{20}
\]

For the zero-rank convention, rejection probability is zero. The
Gaussian assumption, known scale, specified test space, and faithful
event computation are essential to this statement; floating-point
implementation requires separate numerical validation.

For a prespecified significance level \(\alpha\in(0,1)\), reject when
\(p_{\mathrm{selective},j,U}<\alpha\). The procedure tests one uniformly
selected feature for a prespecified target definition. The guarantee is
for that individual test. Reporting multiple discoveries or choosing a
target definition after examining its \(p\)-value requires an additional
error-control argument.

\end{document}
